% year: תשפ"ו
% type: בוחן
% semester: א
% official_solution: yes
% source: exams/בחנים/תשפ״ו, בוחן, 9141.pdf
% part: ב
% number: 7
% points: 10
% categories: continuity, func-limits, multiple-choice

\begin{question}
בדקו את רציפות הפונקציה הבאה בנקודה 3:
\[
f(x)=\begin{cases}
\frac{\sqrt{x+7}-\sqrt{10}}{\sqrt{x+4}-\sqrt{7}}, & \text{if } x>3,\\
\sqrt{x^2+1}, & \text{if } x\le 3
\end{cases}
\]
אילו מהטענות הבאות נכונה?
\begin{enumerate}
  \item הפונקציה רציפה ב-3.
  \item הפונקציה בעלת אי-רציפות סליקה ב-3.
  \item הפונקציה בעלת אי-רציפות מסוג קפיצה ב-3.
  \item הפונקציה בעלת אי-רציפות עיקרית ב-3.
\end{enumerate}
\end{question}

\begin{hint}
בגבול מימין כפלו מונה ומכנה בצמודים של שני הביטויים: $\sqrt{x+7}+\sqrt{10}$ ו-$\sqrt{x+4}+\sqrt7$.
\end{hint}

\begin{solution}
\textbf{התשובה הנכונה: (ג)}.

\textbf{משמאל:} $\sqrt{x^2+1}$ רציפה (הרכבה של פולינום ושורש), ולכן
\[
\lim_{x\to3^-}f(x)=\sqrt{3^2+1}=\sqrt{10}=f(3).
\]
\textbf{מימין:} גבול מהצורה $\frac00$. נכפול בצמודים:
\[
\frac{\sqrt{x+7}-\sqrt{10}}{\sqrt{x+4}-\sqrt7}\cdot\frac{\sqrt{x+4}+\sqrt7}{\sqrt{x+4}+\sqrt7}\cdot\frac{\sqrt{x+7}+\sqrt{10}}{\sqrt{x+7}+\sqrt{10}}
=\frac{\bigl((x+7)-10\bigr)\left(\sqrt{x+4}+\sqrt7\right)}{\bigl((x+4)-7\bigr)\left(\sqrt{x+7}+\sqrt{10}\right)}
=\frac{\sqrt{x+4}+\sqrt7}{\sqrt{x+7}+\sqrt{10}},
\]
כאשר צמצמנו את $x-3\neq0$. הביטוי האחרון רציף ב-3, ולכן
\[
\lim_{x\to3^+}f(x)=\frac{\sqrt7+\sqrt7}{\sqrt{10}+\sqrt{10}}=\frac{2\sqrt7}{2\sqrt{10}}=\frac{\sqrt7}{\sqrt{10}} .
\]
שני הגבולות החד-צדדיים קיימים וסופיים, אך $\frac{\sqrt7}{\sqrt{10}}\neq\sqrt{10}$. לכן הגבול ב-3 אינו קיים, ויש ב-3 אי-רציפות מסוג קפיצה.

\textbf{למה האחרות שגויות:} (א) — הגבול ב-3 אינו קיים, ולכן אין רציפות (יש רק רציפות משמאל). (ב) — אי-רציפות סליקה דורשת קיום הגבול הדו-צדדי. (ד) — אי-רציפות עיקרית דורשת שלפחות אחד הגבולות החד-צדדיים לא יהיה קיים כגבול סופי, וכאן שניהם קיימים.
\end{solution}
